\makeatletter
\def\input@path{{../styles/}{../../styles/}{../../../styles/}{../}{../../}{../../../}}
\makeatother

\documentclass{ee102_pset}
% macros.tex - Course meta information
\renewcommand{\course}{EE 102} % with a space
\renewcommand{\coursetitle}{Signal Processing and Linear Systems}
\renewcommand{\instructor}{Ayush Pandey}
\renewcommand{\student}{Name: }

\renewcommand{\semester}{Fall 2025}
\date{\semester} % this sets the LaTeX date field safely

% Problem set number
\renewcommand{\psetnum}{1}

% Release / due — use renewcommand because package provides empties
\renewcommand{\releasedate}{August 20, 2025}
\renewcommand{\duedate}{August 27, 2025}

\input{preamble.tex}

\author{\rule{3cm}{0.4pt}}
\submitdate{\rule{3cm}{0.4pt}}
\problemset{\shortstack{Problem Set \#13: Practice and Study Guide\\for Midterm Exam \#3}}
\renewcommand{\psetnum}{13}
\renewcommand{\releasedate}{November 25, 2026}
\renewcommand{\duedate}{No submission required}
\shorttitle{Problem Set \#13}
\hypersetup{pdftitle={EE 102 Problem Set 13: Practice and Study Guide for Midterm Exam 3}}

\setlength{\parskip}{8pt plus 1pt minus 1pt}
\setlength{\abovedisplayskip}{6pt plus 2pt minus 2pt}
\setlength{\belowdisplayskip}{6pt plus 2pt minus 2pt}
\setlength{\abovedisplayshortskip}{3pt plus 1pt}
\setlength{\belowdisplayshortskip}{6pt plus 1pt minus 1pt}

\begin{document}
\textbf{Optional practice for Midterm Exam \#3. No submission required.}

\problem{1}

Consider the signal
\[
x(t)=\cos(2\pi\,10t)+10\sin(2\pi t)+3\cos(2\pi\,100t),
\qquad F_s=500\ \mathrm{Hz}.
\]

\problempart[5 points] Write $x_s(t)$, then write the signal as a sequence $x[n]$.

\problempart[10 points] Write $X(e^{j\omega})$ using the DTFT formula from class.

\problempart[10 points] Choose which frequencies you want to keep. Design a discrete-time
filter $H(e^{j\omega})$. Choose your cut-off frequencies, gain, and sketch the
filter frequency response.

\problempart[10 points] Find the output in the frequency domain:
\[
Y(e^{j\omega})=X(e^{j\omega})H(e^{j\omega}).
\]

\problem{2}

Consider the same signal and sampling frequency as in Problem 1.

\problempart[5 points] Find $X(\omega)$.

\problempart[10 points] Write $x_s(t)$ and find $X_s(\omega)$.

\problempart[10 points] Design a continuous-time filter $H(\omega)$ to keep the same
frequencies you chose in Problem 1 and remove the other components and
the extra copies in $X_s(\omega)$. Choose your cut-off frequencies, gain,
and sketch the filter frequency response.

\problempart[10 points] Find the output in the frequency domain:
\[
Y(\omega)=H(\omega)X_s(\omega).
\]

% Adapted from Fall 2025 homework/week12/hw11.tex, Problem 1.
\problem{3}

You know the following properties of a continuous-time signal $x(t)$:
\begin{itemize}
    \item It is dominated by five main frequency components: 100 Hz, 300 Hz,
    600 Hz, 900 Hz, and 1200 Hz, with other frequencies between 100 Hz and
    1500 Hz having negligible amplitudes. There are no frequency components
    above 1500 Hz.
    \item The maximum amplitude of the signal in time-domain is 2.
    \item It is known that the signal can be uniquely reconstructed from its
    samples with a sampling frequency of 10,000 Hz.
\end{itemize}

\problempart[5 points] Sketch $|X(\omega)|$ using the information provided.

\problempart[10 points] Suggest a sampling frequency $f_s<10{,}000$ Hz for this signal
so that it can be perfectly reconstructed from its samples. Justify your answer.

\problempart[15 points] Choose a lower sampling frequency that causes aliasing between
two of the listed frequency components. Describe which components overlap
in the sampled spectrum and why they become indistinguishable after sampling.

\end{document}
